% ---
% id: transmitter/lem:matched-pinching
% kind: lem
% title: "Matched pinching of a block certificate"
% summary: >-
%   A block certificate is unchanged by pinching Alice's sectors or the virtual
%   signal's sectors paired with them, so every state an admitted attack
%   reaches from it is block diagonal with Alice's declared marginal. An
%   operator merely constrained to have that marginal need not be block
%   diagonal.
% prominence: main
% status: accepted
% uses: [def:srcert, def:source-spec]
% mentions: [thm:certificate-sufficiency]
% assumes: [A-SRC]
% refs: []
% provenance: newly-derived
% lean:
% notes: "A-SRC enters through (SR3) only, as the hypothesis that the attack does not act on A, exactly as in thm:certificate-sufficiency (iii). A consumer's licence to restrict its own feasible set to block-diagonal operators is a theorem of that consumer's module and is not stated here."
% ---

A block certificate declares its sectors on Alice's side and on the virtual
signal at once. The lemma shows that this matching survives every admitted
attack, for the states the attack actually reaches; it is about those states
only, and part (iii) shows why it does not restrict any optimisation's
feasible set.

\begin{lemma}\label{lem:matched-pinching}
Let a source specification carry a certificate in block form
(Definition~\ref{def:srcert}), with sector projectors $\Pi^{A}_t$ on $A$ and
$\Pi^{\mathrm v}_t$ on $\Avirt$, $t\in\mathbb{T}$, which sum to the identity
on $A$ and on $\Avirt$ respectively (Definitions~\ref{def:source-spec}
and~\ref{def:srcert}), and write
\begin{equation*}
  \mathrm{Pin}_{A}(\xi)\coloneqq\sum_{t\in\mathbb{T}}\Pi^{A}_t\,\xi\,\Pi^{A}_t ,
  \qquad
  \mathrm{Pin}_{\mathrm v}(\zeta)\coloneqq
  \sum_{t\in\mathbb{T}}\Pi^{\mathrm v}_t\,\zeta\,\Pi^{\mathrm v}_t
\end{equation*}
for the two pinching channels.
\begin{enumerate}[label=(\roman*),leftmargin=2.2em]
\item \emph{(matched structure)} for every $t\in\mathbb{T}$,
\begin{equation}\label{eq:matched}
  (\Pi^{A}_t\otimes\idon{\Avirt})\,\Psi
  =\Psi\,(\Pi^{A}_t\otimes\idon{\Avirt})
  =(\idon{A}\otimes\Pi^{\mathrm v}_t)\,\Psi
  =\Psi\,(\idon{A}\otimes\Pi^{\mathrm v}_t)
  =\Psi_t ,
\end{equation}
the block $\Psi_t$ read as an operator on $\Hs{A\Avirt}$ extended by zero.
Hence
$(\mathrm{Pin}_{A}\otimes\idmap_{\Avirt})(\Psi)=\Psi
=(\idmap_{A}\otimes\mathrm{Pin}_{\mathrm v})(\Psi)$.
\item \emph{(reached states)} assume (SR3), and let $\mathcal{E}$ be any
channel from $\Avirt E$ to $E'$ and $\sigma$ any state on $E$, as in
Theorem~\ref{thm:certificate-sufficiency}\,(iii). The state
$\omega\coloneqq(\idmap_{A}\otimes\mathcal{E})(\Psi\otimes\sigma)$ on $AE'$
satisfies
\begin{equation}\label{eq:reached}
  \tr_{E'}\omega=\tau_A ,\qquad
  (\mathrm{Pin}_{A}\otimes\idmap_{E'})(\omega)=\omega ,\qquad
  \omega=(\idmap_{A}\otimes\mathcal{E}')(\Psi\otimes\sigma)
  \ \text{ with }\
  \mathcal{E}'\coloneqq\mathcal{E}\circ(\mathrm{Pin}_{\mathrm v}\otimes\idmap_{E}) .
\end{equation}
So $\omega=\bigoplus_t\omega_t$ with $\omega_t\in\Pos{A_tE'}$ and
$\tr_{E'}\omega_t=\tau^{(t)}_A$; and the attack $\mathcal{E}$ and the attack
$\mathcal{E}'$ that first measures the sector of the virtual signal reach
the same state.
\item \emph{(operators that are not reached)} the block diagonality of (ii)
does not extend to operators merely constrained to have the marginal: in
general a positive operator on $AE'$ with marginal
$\tau_A=\bigoplus_t\tau^{(t)}_A$ need not be block diagonal along
$\mathbb{T}$, and this happens already for two one-dimensional sectors,
$\tau_A=\tfrac12\idon{A}$ and $E'$ a qubit.
\end{enumerate}
\end{lemma}

\begin{proof}
(i) By Definition~\ref{def:srcert}, $\Psi=\bigoplus_t\Psi_t$ on the \tref{matched
subspace} $\bigoplus_t\Hs{A_t}\otimes\Hs{S_t}$ and zero on its orthogonal
complement. The projector $\Pi^{A}_t\otimes\idon{\Avirt}$ has range
$\bigoplus_{t'}\Hs{A_t}\otimes\Hs{S_{t'}}$, whose intersection with the
matched subspace is $\Hs{A_t}\otimes\Hs{S_t}$, and it commutes with the
projector onto the matched subspace, both being diagonal in the product of
the declared bases. So $(\Pi^{A}_t\otimes\idon{})\Psi$ is the part of
$\Psi$ with range in $\Hs{A_t}\otimes\Hs{S_t}$, which is $\Psi_t$; taking
adjoints, $\Psi$ and $\Psi_t$ being Hermitian, gives the second equality,
and the same argument with $\idon{A}\otimes\Pi^{\mathrm v}_t$, whose range
meets the matched subspace in the same summand, gives the third and fourth.
Summing $(\Pi^{A}_t\otimes\idon{})\Psi(\Pi^{A}_t\otimes\idon{})
=\Psi_t(\Pi^{A}_t\otimes\idon{})=\Psi_t$ over $t$ gives
$(\mathrm{Pin}_{A}\otimes\idmap)(\Psi)=\Psi$, and likewise on the other
factor.

(ii) $\mathcal{E}$ is trace preserving and does not act on $A$, the latter
being (SR3); so $\tr_{E'}\omega=\tr_{\Avirt E}[\Psi\otimes\sigma]
=\tr_{\Avirt}\Psi=\tau_A$ by (SR1). The maps $\mathrm{Pin}_{A}$ and
$\mathcal{E}$ act on complementary tensor factors, so they commute, and by
(i)
\begin{equation*}
  (\mathrm{Pin}_{A}\otimes\idmap_{E'})(\omega)
  =(\idmap_{A}\otimes\mathcal{E})
   \big((\mathrm{Pin}_{A}\otimes\idmap_{\Avirt})(\Psi)\otimes\sigma\big)
  =(\idmap_{A}\otimes\mathcal{E})(\Psi\otimes\sigma)=\omega .
\end{equation*}
Replacing $(\mathrm{Pin}_{A}\otimes\idmap)(\Psi)$ by
$(\idmap\otimes\mathrm{Pin}_{\mathrm v})(\Psi)$, equal to it by (i), in the
middle expression gives
$\omega=(\idmap_{A}\otimes\mathcal{E}')(\Psi\otimes\sigma)$. An operator
fixed by $\mathrm{Pin}_{A}\otimes\idmap_{E'}$ is the direct sum of its
compressions $\omega_t=(\Pi^{A}_t\otimes\idon{})\omega(\Pi^{A}_t\otimes\idon{})$,
each positive semidefinite as a compression of one, and
$\tr_{E'}\omega_t=\Pi^{A}_t\tau_A\Pi^{A}_t=\tau^{(t)}_A$.

(iii) With two one-dimensional sectors spanned by $\ket{0}$ and $\ket{1}$,
$\tau_A=\tfrac12\idon{A}$ and $E'$ a qubit, the projector $\omega$ onto
$(\ket{0}\ket{0}+\ket{1}\ket{1})/\sqrt2$ is positive semidefinite with
marginal $\tau_A$ on $A$, and its compression
$(\Pi^{A}_0\otimes\idon{})\,\omega\,(\Pi^{A}_1\otimes\idon{})
=\tfrac12\ketbra{0}{1}\otimes\ketbra{0}{1}$ is not zero.
\end{proof}
